\claude{
\section{Introduction}\label{sec:intro}
\cx{Can a scalar log reveal when a recurrent network's activity becomes periodic or when a classifier's updates concentrate in fewer directions? A loss curve alone can hide substantial changes in the weights. Two classifiers can start from nearly identical weights and drift apart while both training losses fall smoothly \citep{danovski2024stability}.}

\cx{In grokking, test accuracy rises long after the network fits its training data. \citet{thilak2022slingshot} observed loss spikes as rapid growth in the final layer's weight norm gave way to a plateau. \citet{notsawo2024predicting} used early loss oscillations to anticipate grokking on modular arithmetic. They also estimated dimension with the pooled MacKay--Ghahramani (MG) estimator. Each point was the activation vector of an entire layer for one input.}

\cx{\citet{sussillo2009generating} observed single neurons changing from irregular to periodic activity as a recurrent generator learned. They assessed dimension from principal components of the population activity. We apply MG to points whose coordinates are values of the same scalar quantity at different times \citep{levina2004maximum,mackay2005comments}. For a suitable observation, these delay coordinates can preserve the local dimension of recurrent motion \citep{takens1981detecting,sauer1991embedology,grassberger1983measuring}. During training, the weights themselves change. We ask which changes in the network are visible in a scalar log and whether MG can warn of a loss spike or help choose a model.}
}
